% Part: incompleteness
% Chapter: representability-in-q
% Section: sigma1-completeness

\documentclass[../../../include/open-logic-section]{subfiles}

\begin{document}

\olfileid{inc}{inp}{s1c}
\olsection{تمام \texorpdfstring{$\Sigma_\OLMathNumeral{1}$}{Sigma-1}}

مع أن $\Th{Q}$ وامتداداتها المتسقة القابلة للمحورة غير تامة،
قد رأينا $\Th{Q}$ تثبت كثيرًا من الحقائق الأساسية عن أسماء الأعداد.
ويمكن توسيع هذه النتيجة توسعًا كبيرًا. ولتبيين ما يصل إليه الإثبات في
$\Th{Q}$، نميّز !!{formula}s من الأصناف $\OLId{greek-variable}{Delta}_\OLMathNumeral{0}$ و
$\Sigma_\OLMathNumeral{1}$ و$\Pi_\OLMathNumeral{1}$.
فالـ!!{formula} من الصنف $\Sigma_\OLMathNumeral{1}$، على وجه التقريب،
هي ما صورتها $\lexists[\OLId{latin-ordinary}{x}][!B(\OLId{latin-ordinary}{x})]$، وتكون $!B$ مبنية بالروابط القضوية
والمكممات المحدودة وحدها. وسنثبت أن كل $!A$ تكون !!{sentence}
من الصنف $\Sigma_\OLMathNumeral{1}$ وتصدق في $\Struct{N}$،
تستوفي $\Th{Q}\Proves !A$ (\olref{thm:sigma1-completeness}).

\begin{defn}
\ollabel{defn:bd-quant}
الـ\emph{!!{formula} الوجودية المحدودة} ما كانت على الصورة
$\lexists[\OLId{latin-ordinary}{x}][(\OLId{latin-ordinary}{x}<\OLId{latin-ordinary}{t}\land !A(\OLId{latin-ordinary}{x}))]$، حيث $\OLId{latin-ordinary}{t}$ حد لا يرد فيه المتغير المكمَّم،
واختصار كتابتها المتعارف عليه $\bexists{\OLId{latin-ordinary}{x}<\OLId{latin-ordinary}{t}}{!A(\OLId{latin-ordinary}{x})}$.
%
والـ\emph{!!{formula} الكلية المحدودة} ما كانت على الصورة
$\lforall[\OLId{latin-ordinary}{x}][(\OLId{latin-ordinary}{x}<\OLId{latin-ordinary}{t}\lif !A(\OLId{latin-ordinary}{x}))]$، حيث $\OLId{latin-ordinary}{t}$ حد لا يرد فيه المتغير المكمَّم،
واختصارها المتعارف عليه $\bforall{\OLId{latin-ordinary}{x}<\OLId{latin-ordinary}{t}}{!A(\OLId{latin-ordinary}{x})}$.
\end{defn}

\begin{defn}
\ollabel{defn:delta0-sigma1-pi1-frm}
نعدّ الـ!!{formula}~$!B$ من الصنف $\OLId{greek-variable}{Delta}_\OLMathNumeral{0}$ متى كان بناؤها من
!!{formula}s ذرية، ولم يستعمل فيه إلا الروابط القضوية والتكميم المحدود.
%
ونعدّ الـ!!{formula}~$!A$ من الصنف $\Sigma_\OLMathNumeral{1}$ متى كانت
$!A\ident\lexists[\OLId{latin-ordinary}{x}][!B(\OLId{latin-ordinary}{x})]$، وكانت $!B$ من الصنف $\OLId{greek-variable}{Delta}_\OLMathNumeral{0}$.
%
ونعدّ الـ!!{formula}~$!A$ من الصنف $\Pi_\OLMathNumeral{1}$ متى كانت
$!A\ident\lforall[\OLId{latin-ordinary}{x}][!B(\OLId{latin-ordinary}{x})]$، وكانت $!B$ من الصنف $\OLId{greek-variable}{Delta}_\OLMathNumeral{0}$.
\end{defn}

\begin{lem}
\ollabel{lem:q-proves-clterm-id} إذا كان $\OLId{latin-ordinary}{t}$ حدًا مغلقًا وقيمته
$\Value{\OLId{latin-ordinary}{t}}{\OLId{latin-looped}{N}}=\OLId{latin-ordinary}{n}$، لزم $\Th{Q}\Proves\eq[\OLId{latin-ordinary}{t}][\num \OLId{latin-ordinary}{n}]$.
\end{lem}

\begin{proof}
البرهان بالاستقراء على تعقيد~$\OLId{latin-ordinary}{t}$. ففي الأساس
$\Value{\Obj 0}{\OLId{latin-looped}{N}}=\OLMathNumeral{0}$، وتثبت $\Th{Q}\Proves\eq[\Obj 0][\num \OLMathNumeral{0}]$،
إذ $\num \OLMathNumeral{0}\ident\Obj 0$.
%
وللخطوة نأخذ الحدين $\OLId{latin-ordinary}{t}_\OLMathNumeral{1}$ و$\OLId{latin-ordinary}{t}_\OLMathNumeral{2}$ مع
$\Value{\OLId{latin-ordinary}{t}_\OLMathNumeral{1}}{\OLId{latin-looped}{N}}=\OLId{latin-ordinary}{n}_\OLMathNumeral{1}$ و$\Value{\OLId{latin-ordinary}{t}_\OLMathNumeral{2}}{\OLId{latin-looped}{N}}=\OLId{latin-ordinary}{n}_\OLMathNumeral{2}$،
ونفرض $\Th{Q}\Proves\eq[\OLId{latin-ordinary}{t}_\OLMathNumeral{1}][\num{\OLId{latin-ordinary}{n}_\OLMathNumeral{1}}]$
و$\Th{Q}\Proves\eq[\OLId{latin-ordinary}{t}_\OLMathNumeral{2}][\num{\OLId{latin-ordinary}{n}_\OLMathNumeral{2}}]$.

أما الخلف، فـ$\Value{(\OLId{latin-ordinary}{t}_\OLMathNumeral{1}')}{\OLId{latin-looped}{N}}=\OLId{latin-ordinary}{n}_\OLMathNumeral{1}+\OLMathNumeral{1}$.
ونحصل على $\Th{Q}\Proves\eq[\OLId{latin-ordinary}{t}_\OLMathNumeral{1}'][{\num \OLId{latin-ordinary}{n}_\OLMathNumeral{1}}']$
بقواعد الهوية في الرتبة الأولى، مستعملين فرض الاستقراء والـ!!{formula}
$\eq[\num{\OLId{latin-ordinary}{n}_\OLMathNumeral{1}}'][\num{\OLId{latin-ordinary}{n}_\OLMathNumeral{1}}']$.
ومن تعريف أسماء الأعداد ينتج
$\Th{Q}\Proves\eq[\OLId{latin-ordinary}{t}_\OLMathNumeral{1}'][\num{\OLId{latin-ordinary}{n}_\OLMathNumeral{1}+\OLMathNumeral{1}}]$.

وأما الجمع، فقيمة حده
\[
      \Value{(\OLId{latin-ordinary}{t}_\OLMathNumeral{1} + \OLId{latin-ordinary}{t}_\OLMathNumeral{2})}{\OLId{latin-looped}{N}}
    = \Value{\OLId{latin-ordinary}{t}_\OLMathNumeral{1}}{\OLId{latin-looped}{N}} + \Value{\OLId{latin-ordinary}{t}_\OLMathNumeral{2}}{\OLId{latin-looped}{N}}
    = \OLId{latin-ordinary}{n}_\OLMathNumeral{1} + \OLId{latin-ordinary}{n}_\OLMathNumeral{2}.
\]
فمن فرض الاستقراء وقواعد الهوية نحصل أولًا على
$\Th{Q}\Proves\eq[\OLId{latin-ordinary}{t}_\OLMathNumeral{1}+\OLId{latin-ordinary}{t}_\OLMathNumeral{2}][\num{\OLId{latin-ordinary}{n}_\OLMathNumeral{1}}+\OLId{latin-ordinary}{t}_\OLMathNumeral{2}]$،
ثم نحصل على
$\Th{Q}\Proves\eq[\OLId{latin-ordinary}{t}_\OLMathNumeral{1}+\OLId{latin-ordinary}{t}_\OLMathNumeral{2}][\num{\OLId{latin-ordinary}{n}_\OLMathNumeral{1}}+\num{\OLId{latin-ordinary}{n}_\OLMathNumeral{2}}]$
بتطبيق قواعد الهوية مرة ثانية.
وتعطينا \olref[inc][req][bre]{lem:q-proves-add}
$\Th{Q}\Proves
\eq[\num{\OLId{latin-ordinary}{n}_\OLMathNumeral{1}}+\num{\OLId{latin-ordinary}{n}_\OLMathNumeral{2}}][\num{\OLId{latin-ordinary}{n}_\OLMathNumeral{1}+\OLId{latin-ordinary}{n}_\OLMathNumeral{2}}]$؛
فتلزم $\Th{Q}\Proves\eq[\OLId{latin-ordinary}{t}_\OLMathNumeral{1}+\OLId{latin-ordinary}{t}_\OLMathNumeral{2}][\num{\OLId{latin-ordinary}{n}_\OLMathNumeral{1}+\OLId{latin-ordinary}{n}_\OLMathNumeral{2}}]$.

وأما $\times$، فحجته نظير هذه، مع استعمال
\olref[inc][req][bre]{lem:q-proves-mult}.
%
وهذه جميع وجوه بناء الحدود المغلقة في الحساب. فيثبت
$\Th{Q}\Proves\eq[\OLId{latin-ordinary}{t}][\num \OLId{latin-ordinary}{n}]$ لكل حد مغلق~$\OLId{latin-ordinary}{t}$ قيمته
$\Value{\OLId{latin-ordinary}{t}}{\OLId{latin-looped}{N}}=\OLId{latin-ordinary}{n}$.
\end{proof}

\begin{prob}
فصّل حجة $\times$ في
\olref{lem:q-proves-clterm-id}.
\end{prob}

\begin{lem}
\ollabel{lem:atomic-completeness}
ليكن $\OLId{latin-ordinary}{t}_\OLMathNumeral{1}$ و$\OLId{latin-ordinary}{t}_\OLMathNumeral{2}$ حدين مغلقين. فتصدق الأحكام الآتية:
\begin{enumerate}
\item متى كان $\Value{\OLId{latin-ordinary}{t}_\OLMathNumeral{1}}{\OLId{latin-looped}{N}}=\Value{\OLId{latin-ordinary}{t}_\OLMathNumeral{2}}{\OLId{latin-looped}{N}}$،
    لزم $\Th{Q}\Proves\eq[\OLId{latin-ordinary}{t}_\OLMathNumeral{1}][\OLId{latin-ordinary}{t}_\OLMathNumeral{2}]$.
\item متى كان $\Value{\OLId{latin-ordinary}{t}_\OLMathNumeral{1}}{\OLId{latin-looped}{N}}\neq\Value{\OLId{latin-ordinary}{t}_\OLMathNumeral{2}}{\OLId{latin-looped}{N}}$،
    لزم $\Th{Q}\Proves\eq/[\OLId{latin-ordinary}{t}_\OLMathNumeral{1}][\OLId{latin-ordinary}{t}_\OLMathNumeral{2}]$.
\item متى كان $\Value{\OLId{latin-ordinary}{t}_\OLMathNumeral{1}}{\OLId{latin-looped}{N}}<\Value{\OLId{latin-ordinary}{t}_\OLMathNumeral{2}}{\OLId{latin-looped}{N}}$،
    لزم $\Th{Q}\Proves \OLId{latin-ordinary}{t}_\OLMathNumeral{1}<\OLId{latin-ordinary}{t}_\OLMathNumeral{2}$.
\item متى كان $\Value{\OLId{latin-ordinary}{t}_\OLMathNumeral{2}}{\OLId{latin-looped}{N}}\leq\Value{\OLId{latin-ordinary}{t}_\OLMathNumeral{1}}{\OLId{latin-looped}{N}}$،
    لزم $\Th{Q}\Proves\lnot(\OLId{latin-ordinary}{t}_\OLMathNumeral{1}<\OLId{latin-ordinary}{t}_\OLMathNumeral{2})$.
\end{enumerate}
\end{lem}

\begin{proof}
نأخذ الحدين $\OLId{latin-ordinary}{t}_\OLMathNumeral{1}$ و$\OLId{latin-ordinary}{t}_\OLMathNumeral{2}$، ونسمّي قيمتيهما
$\OLId{latin-ordinary}{n}=\Value{\OLId{latin-ordinary}{t}_\OLMathNumeral{1}}{\OLId{latin-looped}{N}}$ و$\OLId{latin-ordinary}{m}=\Value{\OLId{latin-ordinary}{t}_\OLMathNumeral{2}}{\OLId{latin-looped}{N}}$.

أما $!A\ident \OLId{latin-ordinary}{t}_\OLMathNumeral{1}=\OLId{latin-ordinary}{t}_\OLMathNumeral{2}$، فنعلم من
\olref{lem:q-proves-clterm-id} أن
$\Th{Q}\Proves\eq[\OLId{latin-ordinary}{t}_\OLMathNumeral{1}][\num \OLId{latin-ordinary}{n}]$
و$\Th{Q}\Proves\eq[\OLId{latin-ordinary}{t}_\OLMathNumeral{2}][\num \OLId{latin-ordinary}{m}]$.
إذا كان $\OLId{latin-ordinary}{n}=\OLId{latin-ordinary}{m}$ ثبتت $\Th{Q}\Proves\eq[\num \OLId{latin-ordinary}{n}][\num \OLId{latin-ordinary}{m}]$،
فينتج $\Th{Q}\Proves\eq[\OLId{latin-ordinary}{t}_\OLMathNumeral{1}][\OLId{latin-ordinary}{t}_\OLMathNumeral{2}]$ بتعدي الهوية.
وإذا كان $\OLId{latin-ordinary}{n}\neq \OLId{latin-ordinary}{m}$ ثبتت
$\Th{Q}\Proves\eq/[\num \OLId{latin-ordinary}{n}][\num \OLId{latin-ordinary}{m}]$؛
ومن تعدي الهوية أيضًا يلزم $\Th{Q}\Proves\eq/[\OLId{latin-ordinary}{t}_\OLMathNumeral{1}][\OLId{latin-ordinary}{t}_\OLMathNumeral{2}]$.

وأما $!A\ident \OLId{latin-ordinary}{t}_\OLMathNumeral{1}<\OLId{latin-ordinary}{t}_\OLMathNumeral{2}$، فمرجع الحالتين إلى
البديهية~$!Q_\OLMathNumeral{8}$، ونصها
$\OLId{latin-ordinary}{x}<\OLId{latin-ordinary}{y}\liff\lexists[\OLId{latin-ordinary}{z}][\eq[\OLId{latin-ordinary}{z}'+\OLId{latin-ordinary}{x}][\OLId{latin-ordinary}{y}]]$ لكل $\OLId{latin-ordinary}{x},\OLId{latin-ordinary}{y}$.

إن صدقت $\Sat{\OLId{latin-looped}{N}}{\OLId{latin-ordinary}{t}_\OLMathNumeral{1}<\OLId{latin-ordinary}{t}_\OLMathNumeral{2}}$، وُجد $\OLId{latin-ordinary}{k}\in\Nat$
يحقق $\OLId{latin-ordinary}{n}+\OLId{latin-ordinary}{k}+\OLMathNumeral{1}=\OLId{latin-ordinary}{m}$. وباللمة \olref{lem:q-proves-clterm-id} لدينا
$\Th{Q}\Proves\eq[\OLId{latin-ordinary}{t}_\OLMathNumeral{1}][\num \OLId{latin-ordinary}{n}]$
و$\Th{Q}\Proves\eq[\OLId{latin-ordinary}{t}_\OLMathNumeral{2}][\num \OLId{latin-ordinary}{m}]$.
ومن أول هذه اللمة نحصل على
$\Th{Q}\Proves\eq[{\num \OLId{latin-ordinary}{k}}'+\num \OLId{latin-ordinary}{n}][\num \OLId{latin-ordinary}{m}]$.
فتعطي قواعد الهوية
$\Th{Q}\Proves\eq[{\num \OLId{latin-ordinary}{k}}'+\OLId{latin-ordinary}{t}_\OLMathNumeral{1}][\OLId{latin-ordinary}{t}_\OLMathNumeral{2}]$،
ومنها $\Th{Q}\Proves\lexists[\OLId{latin-ordinary}{z}][\eq[\OLId{latin-ordinary}{z}'+\OLId{latin-ordinary}{t}_\OLMathNumeral{1}][\OLId{latin-ordinary}{t}_\OLMathNumeral{2}]]$.
وعندئذ نستعمل اتجاه~$!Q_\OLMathNumeral{8}$ من اليمين إلى اليسار فنثبت
$\Th{Q}\Proves \OLId{latin-ordinary}{t}_\OLMathNumeral{1}<\OLId{latin-ordinary}{t}_\OLMathNumeral{2}$.

وإن كانت $\Sat/{\OLId{latin-looped}{N}}{\OLId{latin-ordinary}{t}_\OLMathNumeral{1}<\OLId{latin-ordinary}{t}_\OLMathNumeral{2}}$، أي $\OLId{latin-ordinary}{m}\leq \OLId{latin-ordinary}{n}$، فالحجة بالخلف.
%
نعمل داخل~$\Th{Q}$ ونفرض $\OLId{latin-ordinary}{t}_\OLMathNumeral{1}<\OLId{latin-ordinary}{t}_\OLMathNumeral{2}$.
بالاتجاه الآخر من~$!Q_\OLMathNumeral{8}$ يوجد~$\OLId{latin-ordinary}{z}$ يحقق $\eq[\OLId{latin-ordinary}{z}'+\OLId{latin-ordinary}{t}_\OLMathNumeral{1}][\OLId{latin-ordinary}{t}_\OLMathNumeral{2}]$.
ولأن $\Th{Q}\Proves\eq[\OLId{latin-ordinary}{t}_\OLMathNumeral{1}][\num \OLId{latin-ordinary}{n}]$
و$\Th{Q}\Proves\eq[\OLId{latin-ordinary}{t}_\OLMathNumeral{2}][\num \OLId{latin-ordinary}{m}]$، ينتج
$\eq[\OLId{latin-ordinary}{z}'+\num \OLId{latin-ordinary}{n}][\num \OLId{latin-ordinary}{m}]$.
%
ثم نستعمل استقراءً خارجيًا على~$\OLId{latin-ordinary}{m}$، ومعه~$!Q_\OLMathNumeral{5}$ وإلغاء الخلفين بقواعد
الهوية والبديهية الأولى، فنحصل على
$\eq[\OLId{latin-ordinary}{z}'+\num{\OLId{latin-ordinary}{n}-\OLId{latin-ordinary}{m}}][\Obj 0]$.
إن كان $\OLId{latin-ordinary}{m}=\OLId{latin-ordinary}{n}$، صارت $\eq[\OLId{latin-ordinary}{z}'][\Obj 0]$، وتناقض~$!Q_\OLMathNumeral{2}$.
وإن كان $\OLId{latin-ordinary}{m}<\OLId{latin-ordinary}{n}$، صارت
$\eq[(\OLId{latin-ordinary}{z}'+\num{\OLId{latin-ordinary}{n}-\OLId{latin-ordinary}{m}-\OLMathNumeral{1}})'][\Obj 0]$ باستعمال~$!Q_\OLMathNumeral{5}$ مرة أخرى،
فتناقض~$!Q_\OLMathNumeral{2}$ أيضًا. وإذن $\Th{Q}\Proves\lnot(\OLId{latin-ordinary}{t}_\OLMathNumeral{1}<\OLId{latin-ordinary}{t}_\OLMathNumeral{2})$.
\end{proof}

\begin{lem}
\ollabel{lem:bounded-quant-equiv}
لتكن $!A$ !!a{formula}، وليكن $\OLId{latin-ordinary}{t}$ حدًا مغلقًا،
ولتكن $\OLId{latin-ordinary}{k}=\Value{\OLId{latin-ordinary}{t}}{\OLId{latin-looped}{N}}$. عندئذ لدينا التكافؤان:
\begin{enumerate}
\item $\Th{Q}\Proves\bforall{\OLId{latin-ordinary}{x}<\OLId{latin-ordinary}{t}}{!A(\OLId{latin-ordinary}{x})}$ إذا وفقط إذا
    $\Th{Q}\Proves !A(\num \OLMathNumeral{0})\land\dots\land !A(\num{\OLId{latin-ordinary}{k}-\OLMathNumeral{1}})$.
\item $\Th{Q}\Proves\bexists{\OLId{latin-ordinary}{x}<\OLId{latin-ordinary}{t}}{!A(\OLId{latin-ordinary}{x})}$ إذا وفقط إذا
    $\Th{Q}\Proves !A(\num \OLMathNumeral{0})\lor\dots\lor !A(\num{\OLId{latin-ordinary}{k}-\OLMathNumeral{1}})$.
\end{enumerate}
\end{lem}

\begin{proof}
نبرهن التكافؤ الخاص بالمكمم الكلي المحدود.
إذا كانت $\Value{\OLId{latin-ordinary}{t}}{\OLId{latin-looped}{N}}=\OLMathNumeral{0}$، ثبت الطرف الأيسر في~$\Th{Q}$،
إذ لا يوجد $\OLId{latin-ordinary}{x}<\num \OLMathNumeral{0}$ بمقتضى
\olref[inc][req][min]{lem:less-zero}.
وأما الطرف الآخر فالاقتران فيه خال، ونجعله $\ltrue$،
وهو ثابت كذلك في~$\Th{Q}$.
%
يبقى أن تكون $\Value{\OLId{latin-ordinary}{t}}{\OLId{latin-looped}{N}}=\OLId{latin-ordinary}{k}+\OLMathNumeral{1}$ لعدد طبيعي~$\OLId{latin-ordinary}{k}$.
وتجيز لنا \olref{lem:q-proves-clterm-id} أن نستبدل بالصيغة المعطاة
!!a{formula} على الصورة $\bforall{\OLId{latin-ordinary}{x}<\num{\OLId{latin-ordinary}{k}+\OLMathNumeral{1}}}{!A(\OLId{latin-ordinary}{x})}$.

في الاتجاه الأول، لنفرض
$\Th{Q}\Proves\bforall{\OLId{latin-ordinary}{x}<\num{\OLId{latin-ordinary}{k}+\OLMathNumeral{1}}}{!A(\OLId{latin-ordinary}{x})}$،
ولنأخذ $\OLId{latin-ordinary}{n}\leq \OLId{latin-ordinary}{k}$.
ولدينا $\Th{Q}\Proves\num \OLId{latin-ordinary}{n}<\num{\OLId{latin-ordinary}{k}+\OLMathNumeral{1}}$
باللمة \olref{lem:atomic-completeness}،
فيستخرج بالمنطق $\Th{Q}\Proves !A(\num \OLId{latin-ordinary}{n})$.
نكرر ذلك $\OLId{latin-ordinary}{k}+\OLMathNumeral{1}$ مرة، لكل $\OLId{latin-ordinary}{n}\leq \OLId{latin-ordinary}{k}$ مرة،
فنحصل على $\Th{Q}\Proves !A(\num \OLMathNumeral{0})\land\dots\land !A(\num \OLId{latin-ordinary}{k})$،
وهو المطلوب.

وفي الاتجاه المقابل، نفرض
$\Th{Q}\Proves !A(\num \OLMathNumeral{0})\land\dots\land !A(\num \OLId{latin-ordinary}{k})$.
ثم نعمل في $\Th{Q}$ على فرض $\OLId{latin-ordinary}{x}<\num{\OLId{latin-ordinary}{k}+\OLMathNumeral{1}}$.
ومن \olref[inc][req][min]{lem:less-nsucc} نحصل على
$\OLId{latin-ordinary}{x}=\num \OLMathNumeral{0}\lor\dots\lor \OLId{latin-ordinary}{x}=\num \OLId{latin-ordinary}{k}$،
فيثبت بالمنطق~$!A(\OLId{latin-ordinary}{x})$ في كل حالة، ثم يثبت التعميم
$\bforall{\OLId{latin-ordinary}{x}<\num{\OLId{latin-ordinary}{k}+\OLMathNumeral{1}}}{!A(\OLId{latin-ordinary}{x})}$.

وأما التكافؤ الخاص بالـ!!{formula}s ذات التكميم الوجودي المحدود
فحجته على هذا المثال.
\end{proof}

\begin{prob}
فصّل حجة الحالة الوجودية من
\olref{lem:bounded-quant-equiv}.
\end{prob}

\begin{lem}
\ollabel{lem:delta0-completeness}
كل $!A$ تكون !!{sentence} من الصنف $\OLId{greek-variable}{Delta}_\OLMathNumeral{0}$ وتصدق في
$\Struct{N}$، تحقق $\Th{Q}\Proves !A$.
\end{lem}

\begin{proof}
نجري الاستقراء على تعقيد بنية الـ!!{formula} من حيث الروابط والمكممات؛
فتعويض أسماء الأعداد لا يزيد هذه البنية.
%
أما الأساس ففي \olref{lem:atomic-completeness}، فننتقل إلى الخطوة.
ونأخذ الروابط القضوية بعد رد الشرطية والتكافؤ إلى النفي والاقتران والفصل.
وليتضح التفصيل، نقسم النفي بحسب بنية الـ!!{formula} الواقعة في نطاقه.

\begin{enumerate}
\item إذا صدقت $(!A\land !B)$ في $\Struct{N}$، صدقت
$!A$ و$!B$ معًا في~$\Struct{N}$.
فيعطي فرض الاستقراء $\Th{Q}\Proves !A$
و$\Th{Q}\Proves !B$، ويعطي المنطق $\Th{Q}\Proves(!A\land !B)$.
%
\item وإذا صدقت $\lnot(!A\land !B)$ في $\Struct{N}$،
صدقت إحدى $\lnot !A$ و$\lnot !B$ في~$\Struct{N}$.
نأخذ الأولى من غير إخلال بالعموم؛
فيكون $\Th{Q}\Proves\lnot !A$ بفرض الاستقراء،
ثم $\Th{Q}\Proves\lnot(!A\land !B)$ بالمنطق.
%
\item وإذا صدقت $(!A\lor !B)$ في $\Struct{N}$،
فإما أن تصدق $!A$ في $\Struct{N}$، وإما أن تصدق $!B$ في $\Struct{N}$.
وليكن الأول، إذ حجة الثاني مثله.
بفرض الاستقراء $\Th{Q}\Proves !A$،
فتلزم $\Th{Q}\Proves(!A\lor !B)$ بالمنطق.
%
\item وإذا صدقت $\lnot(!A\lor !B)$ في $\Struct{N}$،
صدقت $\lnot !A$ و$\lnot !B$ جميعًا في $\Struct{N}$.
فتثبت $\Th{Q}\Proves\lnot !A$
و$\Th{Q}\Proves\lnot !B$ بفرض الاستقراء،
ومنهما $\Th{Q}\Proves\lnot(!A\lor !B)$ بالمنطق.
%
\item وأما صدق $\bforall{\OLId{latin-ordinary}{x}<\OLId{latin-ordinary}{t}}{!A(\OLId{latin-ordinary}{x})}$ في~$\Struct{N}$،
حيث $\OLId{latin-ordinary}{t}$ حد مغلق و$\OLId{latin-ordinary}{k}=\Value{\OLId{latin-ordinary}{t}}{\OLId{latin-looped}{N}}$، فنستعمل فيه فرض الاستقراء والمنطق.
فمتى صدقت $!A(\num \OLId{latin-ordinary}{n})$ في~$\Struct{N}$ لكل
$\OLId{latin-ordinary}{n}<\Value{\OLId{latin-ordinary}{t}}{\OLId{latin-looped}{N}}$، حصل
$\Th{Q}\Proves !A(\num \OLMathNumeral{0})\land\dots\land !A(\num{\OLId{latin-ordinary}{k}-\OLMathNumeral{1}})$.
فتعطينا \olref{lem:bounded-quant-equiv}
$\Th{Q}\Proves\bforall{\OLId{latin-ordinary}{x}<\OLId{latin-ordinary}{t}}{!A(\OLId{latin-ordinary}{x})}$.
%
\item وأما التكميم الوجودي المحدود، حين تكون !!a{sentence} على الصورة
$\bexists{\OLId{latin-ordinary}{x}<\OLId{latin-ordinary}{t}}{!A(\OLId{latin-ordinary}{x})}$، فحجته نظير حجة التكميم الكلي المحدود.
%
\item وإذا صدقت $\lnot\bforall{\OLId{latin-ordinary}{x}<\OLId{latin-ordinary}{t}}{!A(\OLId{latin-ordinary}{x})}$ في
$\Struct{N}$، حيث $\OLId{latin-ordinary}{t}$ حد مغلق، كانت هذه الـ!!{sentence}
مكافئة للـ!!{sentence} $\bexists{\OLId{latin-ordinary}{x}<\OLId{latin-ordinary}{t}}{\lnot !A(\OLId{latin-ordinary}{x})}$،
والتكافؤ قابل للاشتقاق في~$\Th{Q}$.
فنرجع إلى حجة التكميم الوجودي المحدود.
%
\item وكذلك، إذا صدقت $\lnot\bexists{\OLId{latin-ordinary}{x}<\OLId{latin-ordinary}{t}}!A(\OLId{latin-ordinary}{x})$ في
$\Struct{N}$، حيث $\OLId{latin-ordinary}{t}$ حد مغلق، كانت هذه الـ!!{sentence} مكافئة في
$\Th{Q}$ للصيغة $\bforall{\OLId{latin-ordinary}{x}<\OLId{latin-ordinary}{t}}{\lnot !A(\OLId{latin-ordinary}{x})}$.
فنرجع إلى حجة التكميم الكلي المحدود.
%
\item وأخيرًا، إذا صدقت $\lnot !A$ في $\Struct{N}$،
فقد فرغنا من سائر الحالات إلا أن تكون $!A$ ذرية، أو أن تكون
$\lnot !A\ident\lnot\lnot !B$ لـ!!{sentence}~$!B$
من الصنف $\OLId{greek-variable}{Delta}_\OLMathNumeral{0}$.
إن كانت $!A$ ذرية، حصل $\Th{Q}\Proves\lnot !A$
من \olref{lem:atomic-completeness}.
وإن كانت $\lnot !A\ident\lnot\lnot !B$،
فهي تكافئ~$!B$ تكافؤًا يثبته المنطق في~$\Th{Q}$.
والجملة $!B$ صادقة في~$\Struct{N}$، لأن
$\lnot !A$ صادقة في~$\Struct{N}$.
فيعطي فرض الاستقراء، مع إزالة النفي المزدوج ثم إعادته،
$\Th{Q}\Proves\lnot !A$.
\end{enumerate}
\end{proof}

\begin{prob}
فصّل حجة الحالة الوجودية من
\olref{lem:delta0-completeness}.
\end{prob}

\begin{thm}
\ollabel{thm:sigma1-completeness}
كل $!A$ تكون !!{sentence} من الصنف $\Sigma_\OLMathNumeral{1}$ وتصدق في
$\Struct{N}$، تحقق $\Th{Q}\Proves !A$.
\end{thm}

\begin{proof}
لتكن $\lexists{\OLId{latin-ordinary}{x}}!A(\OLId{latin-ordinary}{x})$ !!{sentence} من الصنف $\Sigma_\OLMathNumeral{1}$
صادقة في~$\Struct{N}$. فيوجد عدد طبيعي~$\OLId{latin-ordinary}{n}$ وإسناد متغيرات~$\OLId{latin-ordinary}{s}$
مع $\OLId{latin-ordinary}{s}(\OLId{latin-ordinary}{x})=\OLId{latin-ordinary}{n}$ و$\Sat{\OLId{latin-looped}{N}}{!A(\OLId{latin-ordinary}{x})}[\OLId{latin-ordinary}{s}]$.
وبالخواص المعروفة لعلاقة التحقيق نحصل على
$\Sat{\OLId{latin-looped}{N}}{!A(\num \OLId{latin-ordinary}{n})}$.
لكن $!A(\num \OLId{latin-ordinary}{n})$ !!{formula} مغلقة من الصنف $\OLId{greek-variable}{Delta}_\OLMathNumeral{0}$،
فتعطي \olref{lem:delta0-completeness}
$\Th{Q}\Proves !A(\num \OLId{latin-ordinary}{n})$.
وبالتعميم الوجودي نحصل على
$\Th{Q}\Proves\lexists[\OLId{latin-ordinary}{x}][!A(\OLId{latin-ordinary}{x})]$.
\end{proof}

\end{document}
